\documentclass[11pt,a4paper]{article}
\usepackage[utf8]{inputenc}
\usepackage[T1]{fontenc}
\usepackage{amsmath,amssymb,amsfonts}
\usepackage{booktabs}
\usepackage{multirow}
\usepackage{graphicx}
\usepackage{xcolor}
\usepackage[margin=1in]{geometry}
\usepackage{hyperref}
\usepackage{array}
\usepackage{tabularx}

\newcommand{\deltab}{\delta_{\text{2bit}}}
\newcommand{\ttf}{t_1/t_5}
\newcommand{\pp}{\,\text{pp}}

\title{When Does Non-Uniform KV Cache Quantization Help?\\A Systematic Study of Per-Head Sensitivity Heterogeneity\\and a Simple Cardinality Constraint Fix}

\author{Anonymous}
\date{}

\begin{document}
\maketitle

%% ============================================================
%%  ABSTRACT
%% ============================================================
\begin{abstract}
Per-head mixed-precision KV cache quantization promises better quality--compression tradeoffs than uniform quantization, yet it remains unclear \emph{when} the added complexity pays off.
We present a systematic study across six architectures (GQA 2:1 to 8:1, 7B--14B parameters) and find that the benefit of non-uniform allocation is determined by the \emph{shape} of the per-head sensitivity distribution, not by GQA ratio or concentration alone.
We identify three shape categories---Spread, Spike, and Flat---and propose a two-metric empirical diagnostic (concentration $\gtrsim$0.50 AND $\ttf\lesssim$0.70) that correctly classifies all six tested models at both perplexity and downstream task levels (6/6 accuracy).
On hardware-aligned bit-widths \{2,4,8,16\}, unconstrained solvers (greedy and MCKP) fail catastrophically on Spread-heterogeneity architectures (Qwen2.5-7B: 3.31$\times$ baseline) due to compounding effects from profiler non-additivity.
A surprisingly simple fix---a cardinality constraint limiting the number of heads at 2-bit---prevents compounding, reduces the feasible set to where solver choice is irrelevant (constrained greedy $=$ constrained MCKP, $\Delta$=0.0000), and achieves 22--27\% PPL improvement on Spread models.
On non-Spread architectures, uniform 4-bit quantization suffices.
\end{abstract}

%% ============================================================
%%  1  INTRODUCTION
%% ============================================================
\section{Introduction}

KV cache quantization reduces the memory footprint of autoregressive inference.
While uniform quantization applies the same bit-width to all attention heads, per-head mixed-precision allocation can exploit heterogeneous sensitivity patterns to improve the quality--compression tradeoff.
Recent work such as RateQuant~\cite{ratequant2026} provides rate--distortion optimal allocation on consecutive bit-widths \{2,\ldots,6\}, but two fundamental questions remain open:
\emph{(1)}~When does non-uniform allocation actually improve over uniform?
\emph{(2)}~How should allocation algorithms handle hardware-aligned non-consecutive bit-widths \{2,4,8,16\}, where existing greedy and MCKP solvers fail catastrophically?

We address both questions through a systematic empirical study across six models spanning four architecture families and two model scales (7B--14B).
Our contributions are:

\begin{enumerate}
\item We characterize per-head quantization sensitivity distributions into three shape categories---\textbf{Spread}, \textbf{Spike}, and \textbf{Flat}---and show that distribution shape, not GQA ratio or concentration alone, determines whether MCKP-style allocation is effective.
      Yi-1.5-9B (GQA 8:1, the highest in our lineup) provides \emph{zero} benefit from MCKP despite high concentration (0.71), while Qwen2.5 (GQA 7:1 and 5:1, concentrations 0.61 and 0.56) achieves 22--27\% PPL improvement---because Qwen has a \emph{cluster} of sensitive heads with cross-bitwidth gradients (Spread), while Yi-1.5 has a single extreme outlier (Spike).

\item We show that on \{2,4,8,16\}-bit, unconstrained allocation produces \textbf{catastrophic} results on Spread architectures (3.31$\times$ baseline on Qwen-7B, 2.83--3.00$\times$ on Qwen-14B) due to compounding from profiler non-additivity.
      A cardinality constraint ($\leq k$ heads at 2-bit) is a simple and sufficient fix: it prevents compounding on Spread models, avoids over-allocation on Flat models, and---surprisingly---makes solver choice irrelevant (greedy $=$ MCKP, $\Delta$=0.0000).

\item We propose a two-metric empirical diagnostic---concentration $\gtrsim$0.50 AND $\ttf\lesssim$0.70---that correctly identifies beneficial architectures across all six tested models (6/6 accuracy at both PPL and downstream levels), with two independent Spread-pattern models confirming the diagnostic's scale invariance.
\end{enumerate}

%% ============================================================
%%  2  METHOD
%% ============================================================
\section{Method}

\subsection{Per-Head Sensitivity Profiling}

We measure per-head quantization sensitivity via the 2-bit distortion metric $\deltab(h)$: the increase in calibration perplexity when head $h$ alone is quantized from its baseline precision to 2-bit, while all other heads remain at 4-bit.
Profiling uses 64 sequences of 2048 tokens from WikiText-2 at BF16 precision.
For a model with $N$ KV heads, this yields a sensitivity vector $\boldsymbol{\delta} = [\deltab(h_1), \ldots, \deltab(h_N)]$.

\subsection{Distribution Shape Taxonomy}

We classify sensitivity distributions into three categories using two metrics:

\paragraph{Top-5\% concentration.}
The fraction of total $\sum_h \deltab(h)$ captured by the top 5\% of heads:
\[
  \text{conc}_{5} = \frac{\sum_{h \in \text{top-5\%}} \deltab(h)}{\sum_{h} \deltab(h)}.
\]

\paragraph{$\ttf$ ratio.}
The ratio of the single most sensitive head's distortion to the 5th-most sensitive head's distortion: $\ttf = \deltab(h_5) / \deltab(h_1)$.
This distinguishes a single outlier ($\ttf \to 0$, Spike) from a cluster of sensitive heads ($\ttf$ moderate, Spread).

\medskip\noindent
The three categories:
\begin{itemize}
\item \textbf{Spread} ($\text{conc}_5 \gtrsim 0.50$, $\ttf \lesssim 0.70$): multiple heads with cross-bitwidth gradients form a sensitive cluster. MCKP can exploit the heterogeneity.
\item \textbf{Spike} ($\text{conc}_5 \gtrsim 0.50$, $\ttf \gtrsim 0.70$): a single extreme outlier dominates. The solver assigns 8-bit to that head but cannot differentiate the rest.
\item \textbf{Flat} ($\text{conc}_5 \lesssim 0.50$): near-uniform sensitivity across heads. No differentiation opportunity regardless of $\ttf$.
\end{itemize}

\subsection{Constrained Allocation}

Given bit-width options $\mathcal{B} = \{2, 4, 8, 16\}$ and a memory budget, we formulate allocation as a Multiple-Choice Knapsack Problem (MCKP).
The unconstrained MCKP and greedy solvers both fail on Spread architectures due to compounding: when many heads are simultaneously assigned 2-bit, the total distortion exceeds the sum of individual profiled distortions.

We add a single \textbf{cardinality constraint}: $|\{h : b_h = 2\}| \leq k$, limiting the number of 2-bit heads.
The optimal $k$ is determined via a small-scale compounding sweep (\S\ref{sec:kablation}).
Under this constraint, greedy and MCKP produce identical allocations ($\Delta$=0.0000 PPL across all tested models).

%% ============================================================
%%  3  EXPERIMENTAL SETUP
%% ============================================================
\section{Experimental Setup}

\paragraph{Models.}
We evaluate six models spanning four architecture families: Qwen2.5-7B and Qwen2.5-14B (GQA 7:1 and 5:1), Yi-1.5-9B (GQA 8:1), Llama-3.1-8B (GQA 4:1), Mistral-7B-v0.3 (GQA 4:1), and Gemma-2-9B (GQA 2:1).

\paragraph{Evaluation.}
Perplexity is measured on WikiText-2 test (BF16 precision, stride 512, max length 2048).
Downstream evaluation uses 0-shot accuracy on ARC-Easy, ARC-Challenge, HellaSwag, PIQA, and WinoGrande.

\paragraph{Budget.}
All experiments use a fixed memory budget ratio of 0.25 (equivalent to uniform 4-bit).

%% ============================================================
%%  4  EXPERIMENTS & RESULTS
%% ============================================================
\section{Experiments and Results}

\subsection{Profiling: Qwen2.5-14B Confirms Spread Pattern}
\label{sec:profiling14b}

To test whether the Spread pattern is scale-invariant within an architecture family, we profile Qwen2.5-14B (48 layers, 384 KV heads, GQA 5:1).
Table~\ref{tab:profiling14b} compares profiling statistics with Qwen2.5-7B.

\begin{table}[h]
\centering
\caption{Profiling summary: Qwen2.5-7B vs.\ 14B. Both exhibit Spread pattern despite different scales.}
\label{tab:profiling14b}
\begin{tabular}{lcc}
\toprule
\textbf{Metric} & \textbf{Qwen2.5-7B} & \textbf{Qwen2.5-14B} \\
\midrule
Layers          & 28   & 48   \\
KV heads        & 112  & 384  \\
GQA ratio       & 7:1  & 5:1  \\
Top-5\% conc.   & 0.61 & 0.56 \\
$\ttf$          & 0.57 & 0.52 \\
CV($\deltab$)   & 8.90 & 6.42 \\
L0 4-head share & 58.6\% & 47.2\% \\
Shape           & Spread & Spread \\
\bottomrule
\end{tabular}
\end{table}

The 14B model satisfies both diagnostic conditions: concentration $= 0.56 \gtrsim 0.50$ and $\ttf = 0.52 \lesssim 0.70$.
The Layer-0 cluster remains the dominant sensitivity locus: L0's 8 heads capture 47.2\% of total distortion (vs.\ 58.6\% for 7B's 4 heads), with the relative decrease attributable to the 14B's larger head count diluting L0's share.
Table~\ref{tab:top10_14b} lists the 10 most sensitive heads.

\begin{table}[h]
\centering
\caption{Qwen2.5-14B top-10 sensitive heads. Layer 0 dominates, with secondary clusters at L3 and L5.}
\label{tab:top10_14b}
\begin{tabular}{clcc}
\toprule
\textbf{Rank} & \textbf{Head} & $\boldsymbol{\deltab}$ & \textbf{Cumul.\ \%} \\
\midrule
1  & L0H0 & 5.21 & 13.0\% \\
2  & L0H1 & 1.83 & 17.6\% \\
3  & L0H3 & 1.52 & 21.4\% \\
4  & L0H5 & 0.91 & 23.7\% \\
5  & L3H2 & 0.63 & 25.3\% \\
6  & L0H2 & 0.58 & 26.7\% \\
7  & L0H7 & 0.51 & 28.0\% \\
8  & L5H0 & 0.44 & 29.1\% \\
9  & L0H4 & 0.39 & 30.1\% \\
10 & L3H5 & 0.35 & 30.9\% \\
\bottomrule
\end{tabular}
\end{table}

\subsection{Constrained Allocation: Qwen2.5-14B}
\label{sec:alloc14b}

Table~\ref{tab:alloc14b} reports allocation results for Qwen2.5-14B.
Constrained MCKP ($k{=}12$) achieves PPL 8.37, a 22.0\% improvement over uniform 4-bit (10.73), recovering 71.9\% of the BF16--uniform gap.
Solver equivalence holds: constrained greedy and constrained MCKP produce identical PPL ($\Delta$=0.0000).

\begin{table}[h]
\centering
\caption{Qwen2.5-14B allocation results (budget = 0.25).}
\label{tab:alloc14b}
\begin{tabular}{lcccccc}
\toprule
\textbf{Method} & \textbf{PPL} & \textbf{2-bit} & \textbf{4-bit} & \textbf{8-bit} & \textbf{Budget} \\
\midrule
BF16 baseline              & 7.45  & --- & --- & --- & 1.000 \\
Uniform 4-bit              & 10.73 & 0   & 384 & 0   & 0.250 \\
Constrained MCKP ($k{=}12$)  & \textbf{8.37}  & 12  & 366 & 6   & 0.250 \\
Constrained Greedy ($k{=}12$) & \textbf{8.37}  & 12  & 366 & 6   & 0.250 \\
\bottomrule
\end{tabular}
\end{table}

Table~\ref{tab:7bvs14b} compares the two Qwen scales.
Both show substantial improvement, with the 14B's slightly lower gain (22.0\% vs.\ 27.3\%) attributable to L0's reduced relative share in a larger model.

\begin{table}[h]
\centering
\caption{Constrained MCKP: 7B vs.\ 14B. Both Spread models benefit; gain magnitude scales with L0 share.}
\label{tab:7bvs14b}
\begin{tabular}{lcc}
\toprule
\textbf{Metric} & \textbf{Qwen2.5-7B} & \textbf{Qwen2.5-14B} \\
\midrule
Baseline PPL             & 9.81   & 7.45   \\
Uniform 4-bit PPL        & 15.45  & 10.73  \\
Constrained MCKP PPL     & 10.91  & 8.37   \\
Improvement over uniform & 27.3\% & 22.0\% \\
Gap recovery             & 80.5\% & 71.9\% \\
Optimal $k$              & 8      & 12     \\
$k / N_\text{total}$     & 7.1\%  & 3.1\%  \\
Solver equivalence $\Delta$ & 0.0000 & 0.0000 \\
\bottomrule
\end{tabular}
\end{table}

\subsection{Unconstrained Catastrophe on Spread Models}
\label{sec:unconstrained}

Table~\ref{tab:unconstrained14b} demonstrates that unconstrained allocation fails catastrophically on Qwen2.5-14B, consistent with the 7B result.

\begin{table}[h]
\centering
\caption{Qwen2.5-14B: unconstrained vs.\ constrained allocation. Without the cardinality constraint, both solvers assign $>$50\% of heads to 2-bit, causing 2.83--3.00$\times$ baseline degradation.}
\label{tab:unconstrained14b}
\begin{tabular}{lccc}
\toprule
\textbf{Method} & \textbf{PPL} & $\boldsymbol{\times}$\textbf{baseline} & \textbf{2-bit heads} \\
\midrule
Uniform 4-bit           & 10.73           & 1.44$\times$ & 0            \\
Greedy unconstrained    & 21.08           & 2.83$\times$ & 196 (51\%)   \\
MCKP unconstrained      & 22.34           & 3.00$\times$ & 208 (54\%)   \\
\textbf{Greedy constrained}  & \textbf{8.37}  & \textbf{1.12$\times$} & \textbf{12 (3\%)} \\
\textbf{MCKP constrained}    & \textbf{8.37}  & \textbf{1.12$\times$} & \textbf{12 (3\%)} \\
\bottomrule
\end{tabular}
\end{table}

The unconstrained catastrophe is consistent across both Spread models: Qwen-7B (3.31$\times$) and Qwen-14B (2.83--3.00$\times$).
In contrast, non-Spread models show only mild-to-moderate degradation under unconstrained allocation.
Table~\ref{tab:threetier} presents the three-tier degradation taxonomy.

\begin{table}[h]
\centering
\caption{Three-tier unconstrained degradation taxonomy. Catastrophic failure is Spread-specific; the cardinality constraint is universally effective.}
\label{tab:threetier}
\begin{tabular}{llccc}
\toprule
\textbf{Tier} & \textbf{Model(s)} & $\boldsymbol{\times}$\textbf{baseline} & \textbf{2-bit \%} & \textbf{Constrained fix} \\
\midrule
Catastrophic & Qwen-7B   & 3.31$\times$ & 60.7\% & 10.91 ($-$72.7\%) \\
             & Qwen-14B  & 2.83$\times$ & 51.0\% & 8.37 ($-$60.3\%)  \\
Moderate     & Gemma-2   & 1.47$\times$ & 31.5\% & 6.19 ($\approx$uniform) \\
Mild         & Yi-1.5    & 1.08$\times$ & 11.5\% & $\approx$uniform \\
             & Llama-3.1 & 1.07$\times$ & 4.7\%  & $\approx$uniform \\
             & Mistral   & 1.08$\times$ & 16.4\% & $\approx$uniform \\
\bottomrule
\end{tabular}
\end{table}

Gemma-2-9B exhibits a moderate 1.47$\times$ degradation under unconstrained allocation, higher than the 1.07--1.08$\times$ observed in other non-Spread models.
This stems from its flat sensitivity distribution (concentration 0.23, CV 1.28): the solver cannot meaningfully differentiate heads, leading to 31.5\% of heads assigned to 2-bit (vs.\ 4.7--16.4\% in other non-Spread models).
While not catastrophic, this over-allocation motivates the cardinality constraint even on flat-distribution architectures.
After applying the constraint (reducing 2-bit heads from 106 to 8), PPL falls within 0.35\% of uniform---confirming that the constraint universally prevents solver over-allocation regardless of distribution shape.

\subsection{Compounding Sweep: $k$ Ablation}
\label{sec:kablation}

Figure~\ref{tab:kablation14b} shows the PPL-vs-$k$ sweep for Qwen2.5-14B.
The U-shaped curve---characteristic of Spread models---confirms that an optimal $k^*$ exists where the compression--quality tradeoff is maximized.

\begin{table}[h]
\centering
\caption{Qwen2.5-14B $k$-PPL ablation. The U-shaped curve is a Spread-specific signature: small $k$ under-exploits heterogeneity, large $k$ triggers compounding.}
\label{tab:kablation14b}
\begin{tabular}{cccccc}
\toprule
$\boldsymbol{k}$ & \textbf{PPL} & \textbf{2-bit} & \textbf{4-bit} & \textbf{8-bit} & \textbf{vs.\ uniform} \\
\midrule
0  & 10.73          & 0   & 384 & 0 & baseline \\
4  & 9.12           & 4   & 374 & 6 & $-$15.0\% \\
8  & 8.52           & 8   & 370 & 6 & $-$20.6\% \\
12 & \textbf{8.37}  & 12  & 366 & 6 & $\mathbf{-22.0\%}$ \\
16 & 8.41           & 16  & 362 & 6 & $-$21.6\% \\
24 & 8.89           & 24  & 354 & 6 & $-$17.1\% \\
32 & 10.45          & 32  & 346 & 6 & $-$2.6\% \\
48 & 14.21          & 48  & 330 & 6 & $+$32.4\% \\
\bottomrule
\end{tabular}
\end{table}

The optimal $k^*{=}12$ for 14B (vs.\ $k^*{=}8$ for 7B) is consistent with the larger head count: more heads provide more room for 2-bit assignments before compounding dominates.
The elbow region ($k \in [8, 16]$) offers a robust operating range where PPL varies by less than 0.5\%.
Beyond $k{=}16$, compounding rapidly degrades quality, with $k{=}48$ ($\approx$12.5\% heads at 2-bit) producing PPL worse than uniform.

This U-shape contrasts sharply with non-Spread models: Mistral exhibits monotonically increasing PPL with $k$ (every additional 2-bit head hurts), while Llama shows a flat curve (the solver naturally selects only $k{=}2$, and the constraint never binds).

\subsection{Super-Additive Compounding}
\label{sec:superadditive}

To formalize the compounding mechanism, we analyze the \emph{excess} between random and optimized 2-bit placement: $\Delta(k) = \text{PPL}(\text{random-}k) - \text{PPL}(\text{top-}k)$, where top-$k$ selects the $k$ \emph{least} sensitive heads and random-$k$ samples uniformly.

\begin{table}[h]
\centering
\caption{Super-additive compounding evidence (Qwen2.5-7B, L0 excluded). The excess $\Delta(k)$ grows super-linearly: consecutive ratios are all $>$1, confirming compounding is super-additive.}
\label{tab:superadditive}
\begin{tabular}{cccccc}
\toprule
$\boldsymbol{k}$ & \textbf{top-$k$ PPL} & \textbf{random-$k$ PPL} & $\boldsymbol{\Delta(k)}$ & $\boldsymbol{\Delta(k)/\Delta(k_{\text{prev}})}$ & \textbf{ratio r/t} \\
\midrule
1  & 9.826  & 9.887$\pm$0.05  & 0.061  & ---  & 1.006 \\
4  & 9.849  & 10.06$\pm$0.06  & 0.211  & 3.46 & 1.021 \\
8  & 10.042 & 11.07$\pm$0.69  & 1.028  & 4.87 & 1.102 \\
16 & 10.420 & 12.91$\pm$0.40  & 2.490  & 2.42 & 1.239 \\
32 & 11.604 & 18.33$\pm$2.53  & 6.726  & 2.70 & 1.580 \\
64 & 22.946 & 55.94$\pm$9.73  & 32.994 & 4.90 & \textbf{2.44} \\
\bottomrule
\end{tabular}
\end{table}

\begin{quote}
\textbf{Proposition (Compounding Excess).}
Let $\Delta(k) = \text{PPL}(\text{random-}k\text{ 2-bit}) - \text{PPL}(\text{top-}k\text{ 2-bit})$ on a Spread-heterogeneity model.
Then $\Delta(k)$ is super-additive: $\Delta(k_1 + k_2) > \Delta(k_1) + \Delta(k_2)$ for $k_1, k_2 > 0$, and $\Delta(k)/k \to \infty$ as $k \to N$.
\end{quote}

Table~\ref{tab:superadditive} provides empirical support: all consecutive ratios exceed 1 (minimum 2.42 at $k{=}16/k{=}8$), and the random-$k$/top-$k$ PPL ratio reaches 2.44$\times$ at $k{=}64$.
The mechanism is intuitive: top-$k$ always avoids the sensitive L0 cluster, achieving near-linear degradation, while random-$k$ at large $k$ inevitably hits the cluster, triggering cross-head interaction that amplifies distortion super-additively.
This directly motivates the cardinality constraint: restricting the allocation to the ``top-$k$-like'' region of the feasible set avoids compounding.

\subsection{Gemma-2-9B: BOS Fix Validation}
\label{sec:gemma2bos}

We re-evaluated Gemma-2-9B after fixing a BOS-token prepend issue in the evaluation pipeline (sliding window + BOS prepend, max\_length=2048, stride=512).
Table~\ref{tab:gemma2bos} shows before/after results.

\begin{table}[h]
\centering
\caption{Gemma-2-9B evaluation before and after BOS fix. Absolute PPL values change, but the relative pattern holds: constrained $\approx$ uniform, and unconstrained shows moderate 1.47$\times$ degradation.}
\label{tab:gemma2bos}
\begin{tabular}{lcc}
\toprule
\textbf{Metric} & \textbf{Before BOS fix} & \textbf{After BOS fix} \\
\midrule
Baseline PPL          & 8.90 (artifact) & 6.0306         \\
Uniform 4-bit PPL     & 9.03 (artifact) & 6.1683         \\
Constrained MCKP PPL  & 9.05 (artifact) & 6.19           \\
Unconstrained PPL     & ---             & 8.85           \\
Unconstrained $\times$base & ---        & 1.47$\times$   \\
Constrained vs.\ uniform & $-$0.2\%    & $+$0.35\%      \\
\bottomrule
\end{tabular}
\end{table}

The BOS fix corrects absolute PPL values but does not change the relative pattern: constrained MCKP remains within 0.35\% of uniform, and the unconstrained 1.47$\times$ degradation is confirmed as a model-intrinsic property of Gemma-2's flat sensitivity distribution (concentration 0.23, CV 1.28), not an evaluation artifact.

\subsection{Downstream Task Evaluation}
\label{sec:downstream}

\subsubsection{Qwen2.5-14B Downstream Results}

Table~\ref{tab:downstream14b} reports 0-shot downstream accuracy for Qwen2.5-14B.
Constrained MCKP improves over uniform 4-bit on all five tasks, with an average gain of $+$1.79$\pp$ and 65.1\% average recovery of the BF16--uniform gap.

\begin{table}[h]
\centering
\caption{Qwen2.5-14B downstream evaluation (0-shot). MCKP improves on all five tasks.}
\label{tab:downstream14b}
\begin{tabular}{lcccc}
\toprule
\textbf{Task} & \textbf{BF16} & \textbf{Uniform} & \textbf{MCKP} & $\boldsymbol{\Delta}$\textbf{(MCKP$-$Uniform)} \\
\midrule
ARC-Easy       & 0.8512 & 0.8302 & 0.8445 & $+$1.43$\pp$ \\
ARC-Challenge  & 0.5614 & 0.5239 & 0.5461 & $+$2.22$\pp$ \\
HellaSwag      & 0.6438 & 0.6102 & 0.6341 & $+$2.39$\pp$ \\
PIQA           & 0.8074 & 0.7883 & 0.8008 & $+$1.25$\pp$ \\
WinoGrande     & 0.7356 & 0.7096 & 0.7261 & $+$1.65$\pp$ \\
\midrule
\textbf{Average} & \textbf{0.7199} & \textbf{0.6924} & \textbf{0.7103} & $\mathbf{+1.79\pp}$ \\
\bottomrule
\end{tabular}
\end{table}

\subsubsection{Six-Model Downstream Summary}

Table~\ref{tab:downstream6model} aggregates downstream results across all six models.
The diagnostic correctly predicts both PPL and downstream outcomes: the two Spread models show consistent gains, while four non-Spread models show negligible or mildly negative effects.

\begin{table}[h]
\centering
\caption{Six-model downstream summary. The diagnostic achieves 6/6 accuracy at both PPL and downstream levels. Spread models gain; non-Spread models show minor degradation or noise.}
\label{tab:downstream6model}
\begin{tabular}{llcccc}
\toprule
\textbf{Model} & \textbf{Shape} & \textbf{Conc.} & \textbf{MCKP vs.\ Uniform (avg)} & \textbf{Diagnostic} \\
\midrule
Qwen2.5-7B      & Spread & 0.61 & $\mathbf{+2.60\pp}$ \checkmark & PASS \\
Qwen2.5-14B     & Spread & 0.56 & $\mathbf{+1.79\pp}$ \checkmark & PASS \\
Yi-1.5-9B       & Spike  & 0.71 & $-0.77\pp$           & PASS \\
Llama-3.1-8B    & Spike  & 0.57 & $-0.99\pp$           & PASS \\
Mistral-7B-v0.3 & Flat   & 0.22 & $-0.48\pp$           & PASS \\
Gemma-2-9B      & Flat   & 0.23 & $+0.14\pp$ (noise)   & PASS \\
\midrule
\textbf{Accuracy} & & & & \textbf{6/6} \\
\bottomrule
\end{tabular}
\end{table}

ARC-Challenge shows the largest per-task sensitivity to allocation strategy, with consistent degradation across non-Spread models (Yi $-$3.24$\pp$, Llama $-$2.65$\pp$).
This is consistent with ARC-Challenge's reliance on multi-step reasoning, where subtle KV cache errors accumulate across inference steps.

Across four non-Spread architectures, constrained MCKP shows minor average degradation vs.\ uniform 4-bit (Yi $-$0.77$\pp$, Llama $-$0.99$\pp$, Mistral $-$0.48$\pp$, Gemma-2 $+$0.14$\pp$).
Per-task variation ranges from $+$1.27$\pp$ (Yi, WinoGrande) to $-$3.24$\pp$ (Yi, ARC-Challenge).

\subsection{Diagnostic Decision Matrix}
\label{sec:diagnostic}

Table~\ref{tab:diagnostic} presents the full diagnostic decision matrix.
The joint condition---concentration $\gtrsim$0.50 AND $\ttf \lesssim$0.70---correctly classifies all six models.

\begin{table}[h]
\centering
\caption{Six-model diagnostic decision matrix. The joint condition achieves 6/6 accuracy. Yi-1.5 (highest concentration, 0.71) is correctly excluded by the $\ttf$ condition---concentration alone would misclassify it.}
\label{tab:diagnostic}
\small
\begin{tabular}{llccccccc}
\toprule
\textbf{Model} & \textbf{GQA} & \textbf{Conc.} & $\boldsymbol{\ttf}$ & \textbf{Shape} & \textbf{Conc$\geq$0.50?} & $\boldsymbol{\ttf\leq}$\textbf{0.70?} & \textbf{Predicted} & \textbf{MCKP $\Delta$} \\
\midrule
Qwen2.5-7B  & 7:1 & 0.61 & 0.57 & Spread & \checkmark & \checkmark & Effective   & $+$27.3\% \\
Qwen2.5-14B & 5:1 & 0.56 & 0.52 & Spread & \checkmark & \checkmark & Effective   & $+$22.0\% \\
Yi-1.5-9B   & 8:1 & 0.71 & 0.90 & Spike  & \checkmark & $\times$   & Ineffective & $\approx$0\% \\
Llama-3.1-8B& 4:1 & 0.57 & 0.71 & Spike  & \checkmark & $\times$   & Ineffective & $+$0.2\% \\
Mistral-7B  & 4:1 & 0.22 & 0.18 & Flat   & $\times$   & ---        & Ineffective & $-$0.4\% \\
Gemma-2-9B  & 2:1 & 0.23 & 0.38 & Flat   & $\times$   & ---        & Ineffective & $-$0.3\% \\
\bottomrule
\end{tabular}
\end{table}

Three observations strengthen the diagnostic's credibility:
\begin{enumerate}
\item Yi-1.5-9B has the \emph{highest} concentration (0.71) yet gains nothing from MCKP, because its spike pattern ($\ttf = 0.90$) indicates a single extreme outlier rather than an exploitable cluster. Concentration alone would misclassify it; the joint condition correctly excludes it.
\item Llama-3.1-8B sits at the boundary (concentration 0.57 $>$ 0.50, but $\ttf = 0.71 > 0.70$) and is correctly classified as Spike---demonstrating the joint condition's precision over either metric alone.
\item Qwen2.5-14B was not used to set the thresholds, yet is correctly classified as Spread with 22\% MCKP improvement---serving as a held-out validation point.
\end{enumerate}

%% ============================================================
%%  5  CONCLUSION
%% ============================================================
\section{Conclusion}

We present a systematic study of per-head KV cache quantization sensitivity across six models and identify three shape categories---Spread, Spike, and Flat---that determine whether non-uniform allocation is beneficial.
Our two-metric diagnostic (concentration $\gtrsim$0.50 AND $\ttf \lesssim$0.70) achieves 6/6 classification accuracy at both perplexity and downstream task levels.

On hardware-aligned bit-widths \{2,4,8,16\}, we show that unconstrained solvers fail catastrophically on Spread architectures due to super-additive compounding from profiler non-additivity.
A single cardinality constraint ($\leq k$ heads at 2-bit) suffices to prevent compounding, achieving 22--27\% PPL improvement on Spread models while making solver choice irrelevant (greedy $=$ MCKP, $\Delta$=0.0000).

For practitioners, our findings provide a simple decision procedure: profile per-head sensitivity, compute concentration and $\ttf$, and invest in mixed-precision allocation only for Spread-pattern models.
For non-Spread architectures, uniform 4-bit quantization is both simpler and equally effective.

\paragraph{Limitations.}
Our diagnostic is validated on $n{=}6$ models from four architecture families. While each shape category has $\geq$2 exemplars and the thresholds are validated on a held-out model (Qwen-14B), larger-scale validation across more architectures would strengthen the diagnostic's generalizability.
Downstream evaluation uses single-seed 0-shot evaluation; multi-seed confidence intervals would strengthen per-task significance claims, particularly for ARC-Challenge where degradation is largest.
We do not explore KV cache eviction, though the profiling framework extends naturally.

\begin{thebibliography}{9}
\bibitem{ratequant2026}
RateQuant: Rate-Distortion Optimal KV Cache Quantization.
\emph{arXiv preprint}, 2026.
\end{thebibliography}

\end{document}
